\section{Introduction} \label{introduction}

The research on a more efficient operation of vessel is a direction that is being actively pursued in the maritime world. 
Efficient ship operation translates to a noticeable increase in profitability and sustainability. 
The Fuel Oil Consumption (FOC) is a determining factor in the ship's operating cost, with fuel cost making up to 75\% of the total ship operating cost \cite{Abebe.2020,Gkerekos.2019}. 
Different ship parameters affect the ship's fuel consumption and research shows that the reduction of ship speed is the most economical method to reduce fuel consumption. 
Precise forecasting and modelling of the ship speed will lead to a significant reduction in a ship's operating cost and could possibly extend the longevity of the ship. \\

In this data-driven thesis, powerful machine learning method will be utilized for modelling and forecasting of the ship's speed. To develop the model. 
The data used will be from the voyage of a Hammerhus RoPax ship between the port of Koege and Roenne. 
The modelling and forecasting using Random Forest Regressor will be emphasis of this thesis. The developed model of the ship's speed is then to be validated and where applicable, 
further optimized to increase its performance. 
This model is then to be used to predict the ship's fuel consumption.



%\begin{figure}[h]
%    \centering
%    \includegraphics[width=0.50\textwidth]{02_figures/logo_en.jpg}
%    \caption{A nice plot}
%    \label{fig:mesh1}
%\end{figure}

